% !PRINTED 477
\weberpage{477}

%<<CONT
und nach \S.\,146, 1. ist also für jeden Index $k$
\weq{2}{\varrho_k = \Phi[\chi(\varrho_k), \chi_1(\varrho_k) \ldots \chi_{m-1}(\varrho_k)].}
Wenn nun die Substitutionen $\sigma_h, \sigma_k$ dieselbe Permutation unter den $\alpha$ hervorrufen, so ist für jeden Index $s = 0, 1 \ldots m-1$
\[
\chi_s(\varrho_h) = \chi_s(\varrho_k),
\]
und daraus folgt nach (2), dass auch $\sigma_h = \sigma_k$ sein muss.

Wird durch die Substitution $\sigma_h$ die Permutation $\pi_h$ unter den Wurzeln von $F(x) = 0$ hervorgerufen, so sagen wir, dass \gesperrt{die Permutation $\pi_h$ der Substitution $\sigma_h$ entspricht}.

Jeder Substitution entspricht eine bestimmte Permutation, und die Anzahl dieser Permutationen ist also, die identische Permutation mitgerechnet, die der identischen Substitution entspricht, gleich $\mu$.

Wenn den Substitutionen $\sigma_h$ und $\sigma_k$ die Permutationen $\pi_h$ und $\pi_k$ entsprechen, so entspricht der zusammengesetzten Substitution $\sigma_h \sigma_k$ die zusammengesetzte Permutation $\pi_h \pi_k$. Denn die Permutationen sind das Ergebniss der entsprechenden Substitutionen auf gewisse Grössen des Körpers $\Omega(\varrho)$, und dabei ist es gleichgültig, ob wir zwei Substitutionen nach einander oder mit einem Male die aus beiden zusammengesetzte Substitution ausführen (nach \S.\,147). Wir haben also den Satz:

\gesperrt{I. Es entsprechen den $\mu$ Substitutionen des Körpers $\Omega(\varrho)$}
\weq{3}{\sigma, \sigma_1, \sigma_2 \ldots \sigma_{\mu-1},}
\gesperrt{$\mu$ Permutationen von $m$ Ziffern}
\weq{4}{\pi, \pi_1, \pi_2 \ldots \pi_{\mu-1}}
\gesperrt{in der Weise, dass der aus zwei Substitutionen zusammengesetzten Substitution die aus den entsprechenden Permutationen zusammengesetzte Permutation entspricht.}

Daraus folgt aber, dass die Permutationen (4) eine Gruppe bilden. Wir bezeichnen diese Gruppe mit $P$ und nennen sie die \gesperrt{Galois'sche Gruppe der Gleichung $F(x) = 0$, oder auch die Galois'sche Gruppe eines jeden der Körper $\Omega(\alpha), \Omega(\alpha_1) \ldots \Omega(\alpha_{m-1})$}.

Diese Gruppe steht zu der Gruppe der Substitutionen des Körpers $\Omega(\varrho)$ in der durch den Satz I. ausgedrückten Beziehung, weshalb die beiden Gruppen \gesperrt{isomorph} genannt werden.
